% ===== ch08 : 原文 main.tex 第 2170--2670 行（5b 特征群作用 / 5c 权空间分解）=====
% 本文件是原文 LaTeX 源码的逐行分片（去掉整段注释掉的旧稿），仅供翻译时对照；最终译文写到 src/chunks/ch08.tex。
\subsection{The action of character group $\widehat{G}$ on $\Mod_G(X)$.} In the abelian category $\Mod_G (X)$ of $G$-equivariant sheaves on $X$ there is a canonical action of a specific abelian group, $\widehat{G}$, called the character group of $G$, which we will define as follows.
\begin{defn}
A \emph{character} of a group scheme $(G,m,e)$ is a group scheme homomorphism
$
\cha : G \longrightarrow \Gm
$
from $G$ to the multiplicative group scheme
$\Gm := \Spec\left(\field[t, t^{-1}]\right)$.
Equivalently,
%\footnote{
% ---- 原文行 2180 ----
%By the universal property of
%$\Gm = \Spec\left(\field[t, t^{-1}]\right)$,
%a morphism of schemes
%$\cha : G \rightarrow \Gm$
%corresponds bijectively to a $\field$-algebra homomorphism $\phi : \field[t, t^{-1}] \rightarrow \Gamma(G, \mathcal{O}_G)$,
%which is uniquely determined by the choice of an invertible section
%$\cha = \phi(t) \in \Gamma\left(G, \mathcal{O}_G^\times\right)$.
%Under this identification,
%the group scheme homomorphism condition
%$(\cha \times \cha)\cdot m_{\Gm} = m\cdot \cha$
%is equivalent to the condition
%$m^* \cha = \left(\pr_{G,1}^{G \times G}\right)^* \cha \cdot \left(\pr_{G,2}^{G \times G}\right)^* \cha$.
%},
a character corresponds to a global section
$
\cha
\in
\Gamma\left(G, \mathcal{O}_G^\times\right)
$
of the trivial $\Gm$-bundle $\mathcal{O}_X^\times$ on $X$
% ---- 原文行 2200 ----
satisfying
$$
m^* \cha
=
\left(\pr_{G,1}^{G \times G}\right)^* \cha
\cdot
\left(\pr_{G,2}^{G \times G}\right)^* \cha
\in
\Gamma\left(G \times G, \mathcal{O}_{G \times G}^\times\right),
$$
where
$\pr_{G,1}^{G \times G}$,
$\pr_{G,2}^{G \times G}
:
G \times G \rightarrow G
$
are the projections onto the first and second factors, respectively,
$
\left(\pr_{G,i}^{G \times G}\right)^* \cha
\in
% ---- 原文行 2220 ----
\Gamma\left(
G \times G,
\left(\pr_{G,i}^{G \times G}\right)^*\mathcal{O}_G^\times
\cong
\mathcal{O}_{G \times G}^\times
\right)
$
is the pullback of $\cha$
by $\pr_{G,i}^{G\times G}$
($i\in\{1,2\}$),
and $\cdot$ denotes the multiplication of elements of the ring
$\mathcal{O}_{G\times G}^\times(G\times G)$.






The set of all characters of $G$ forms an abelian group,
with the group structure induced by the multiplicative group scheme structure of $\Gm$,
% ---- 原文行 2240 ----
denoted by
$$
\widehat{G}
:=
\operatorname{Hom}_{\mathsf{GrpSch}(\field)}\left(G, \Gm\right),
$$
and is called the \emph{character group} of $G$.

\end{defn}




\begin{defn}[Twisting by a character]


For each character
$\cha \in \widehat{G}\subset\Gamma\left(G, \mathcal{O}_G^\times\right)$,
the \emph{trivial line bundle twisted by $\cha$}
is defined by
% ---- 原文行 2260 ----
$$
\mathcal{O}_X(\cha)
:=
(\mathcal{O}_X, \id_{\mathcal{O}_{G\times X}}(\cha))
\in \vb_G(X),
$$
where the underlying $\mathcal{O}_X$-module is $\mathcal{O}_X$
and the $G$-linearization
$$
\id_{\mathcal{O}_{G\times X}}(\cha)
:=
\left(\pr_G^{G \times X}\right)^* \cha \cdot \id_{\mathcal{O}_{G\times X}}
:
\mathcal{O}_{G\times X}
\xrightarrow{\cong}
\mathcal{O}_{G\times X}
$$
(under the canonical identifications
$
\sigma^* \mathcal{O}_X
% ---- 原文行 2280 ----
\cong
\mathcal{O}_{G\times X}
$
and
$
\left(\pr_X^{G \times X}\right)^* \mathcal{O}_X
\cong
\mathcal{O}_{G\times X}
$)
is defined by
the multiplication of
the pullback
$$
\left(\pr_G^{G \times X}\right)^* \cha
\in
\Gamma\left(G \times X, \left(\pr_G^{G\times X}\right)^*\mathcal{O}_G^\times \cong \mathcal{O}_{G \times X}^\times\right)
$$
of $\cha$ by the canonical projection
$\pr_G^{G\times X}:G\times X\rightarrow G$.
More precisely,
% ---- 原文行 2300 ----
we define
$$
\left(\left(\id_{\mathcal{O}_{G\times X}}(\cha)\right)(U)\right)(s)
:=
\left(\left.\left(\pr_G^{G \times X}\right)^* \cha\right|_U\right)
\cdot
s
\in
\mathcal{O}_{G\times X}(U)
$$
for each open
$U\subseteq G\times X$
and
$s\in\mathcal{O}_{G\times X}(U)$,
where
$\cdot$ denotes the multiplication of elements of the ring
$\mathcal{O}_{G\times X}(U)$.



% ---- 原文行 2320 ----

Let $(\mathcal{F}, \theta)$ be a $G$-equivariant $\mathcal{O}_X$-module in $\Mod_G(X)$ and $\cha \in \widehat{G}$.
The \emph{twisting of $(\mathcal{F},\theta)$ by $\cha$}
is defined by
$$
(\mathcal{F},\theta)(\cha)
:=
(\mathcal{F}, \theta(\cha))
\in \Mod_G(X),
$$
where the underlying $\mathcal{O}_X$-module is $\mathcal{F}$
and the $G$-linearization
$$
\theta(\cha)
:=
\left(\pr_G^{G \times X}\right)^* \cha \cdot \theta
:
\sigma^* \mathcal{F}
\xrightarrow{\cong}
\left(\pr_X^{G \times X}\right)^* \mathcal{F}
% ---- 原文行 2340 ----
$$
is defined by the scalar multiplication of
the pullback
$$
\left(\pr_G^{G \times X}\right)^* \cha
\in
\Gamma\left(G \times X, \left(\pr_G^{G\times X}\right)^*\mathcal{O}_G^\times \cong \mathcal{O}_{G \times X}^\times\right)
$$
%of $\cha$ by the canonical projection
%$\pr_G^{G\times X}:G\times X\rightarrow G$
on the original $G$-linearization $\theta$.
More precisely,
we define
$$
(\theta_{\cha}(U))(s)
:=
\left(\left.\left(\pr_G^{G \times X}\right)^* \cha\right|_U\right)
\cdot
(\theta(U)(s))
\in
% ---- 原文行 2360 ----
\left(\left(\pr_X^{G \times X}\right)^*\mathcal{F}\right)(U)
$$
for each open
$U\subseteq G\times X$
and
$s\in(\sigma^*\mathcal{F})(U)$,
where
$\cdot$ denotes the scalar multiplication of elements of the $\mathcal{O}_{G\times X}(U)$-module
$\left(\left(\pr_X^{G \times X}\right)^*\mathcal{F}\right)(U)$
by elements of the ring $\mathcal{O}_{G\times X}(U)$.










% ---- 原文行 2380 ----


Twisting by characters defines a strict $\widehat{G}$-action on $\Mod_G(X)$
$$
\phi:\widehat{G}\rightarrow\Aut_{\str}(\Mod_G(X)),
\quad
\cha\mapsto \phi_{\cha}
$$
where
$$
\phi_{\cha}:\Mod_G(X)\rightarrow\Mod_G(X),
\quad
(\mathcal{F},\theta)
\mapsto
(\mathcal{F},\theta(\cha)).
$$
This restricts to strict $\widehat{G}$-actions on
$\QCoh_G(X)$, $\Coh_G(X)$, and $\VB_G(X)$.
\end{defn}

% ---- 原文行 2400 ----



\begin{rmk}
There is a natural isomorphism
$$
(\mathcal{F},\theta)(\cha)
\cong
(\mathcal{F},\theta)
\underset{\mathcal{O}_X}{\otimes}
\mathcal{O}_X(\cha)
\in\Mod_G(X).
$$
\end{rmk}


\subsection{Weight space decompositions} Since $\widehat{G}$ is abelian, its action on hom spaces in $\Mod_G (X)$ will induce a decomposition of hom spaces into one-dimensional irreducible $\widehat{G}$-representations, which we call customly \emph{weight space decompositions}.

\begin{defn}
Let $\Lambda$ be an abelian group and $\field$ be a commutative ring.
% ---- 原文行 2420 ----
The \emph{group algebra} $\field[\Lambda]$ is a Hopf algebra with a basis $\{e_\lambda\}_{\lambda \in \Lambda}$,
where the comultiplication $\Delta : \field[\Lambda] \rightarrow \field[\Lambda] \underset{\field}{\otimes} \field[\Lambda]$ and the counit $\varepsilon : \field[\Lambda] \rightarrow \field$ are defined by
$$
\Delta(e_\lambda) = e_\lambda \otimes e_\lambda \quad \text{and} \quad \varepsilon(e_\lambda) = 1
\quad\textnormal{for all $\lambda \in \Lambda$}.
$$

A \emph{right comodule} over $\field[\Lambda]$ is a $\field$-module $V$ equipped with a $\field$-linear map $\rho : V \rightarrow V \underset{\field}{\otimes} \field[\Lambda]$ satisfying the coassociativity and counit axioms:
$$
\rho \comp (\rho \otimes \id_{\field[\Lambda]}) = \rho \comp (\id_V \otimes \Delta)
\quad\text{and}\quad
\rho \comp (\id_V \otimes \varepsilon) = \id_V.
$$
\end{defn}

\begin{prop}[Weight space decomposition of $\field[\Lambda\text{$]$}$-comodules]
\label{prop:WeightSpaceDecompositionForComodule}
A $\field$-module $V$ has a right comodule structure over $\field[\Lambda]$ if and only if $V$ is $\Lambda$-graded, i.e., it decomposes into a direct sum of weight spaces
$$
V = \bigoplus_{\lambda \in \Lambda} V_\lambda,
% ---- 原文行 2440 ----
$$
where $V_\lambda = \{v \in V \mid \rho(v) = v \otimes e_\lambda\}$.
\end{prop}

\begin{proof}
For any $v \in V$, its image under $\rho$ can be uniquely written as a finite sum
$$
\rho(v) = \sum_{\lambda \in \Lambda} P_\lambda(v) \otimes e_\lambda,
$$
where $P_\lambda : V \rightarrow V$ are $\field$-linear maps.
Applying the counit axiom $\rho \comp (\id_V \otimes \varepsilon) = \id_V$ to $v$ yields
$$
v = (\id_V \otimes \varepsilon)(\rho(v)) = (\id_V \otimes \varepsilon)\left(\sum_{\lambda \in \Lambda} P_\lambda(v) \otimes e_\lambda\right) = \sum_{\lambda \in \Lambda} P_\lambda(v) \cdot 1 = \sum_{\lambda \in \Lambda} P_\lambda(v).
$$
This shows that the sum of the images of $P_\lambda$ covers $V$.
Next, applying the coassociativity axiom
$
\rho \comp (\rho \otimes \id_{\field[\Lambda]}) = \rho \comp (\id_V \otimes \Delta)
$
to $v$ gives
% ---- 原文行 2460 ----
$$
\sum_{\lambda, \mu \in \Lambda} P_\mu(P_\lambda(v)) \otimes e_\mu \otimes e_\lambda = \sum_{\lambda \in \Lambda} P_\lambda(v) \otimes e_\lambda \otimes e_\lambda.
$$
Comparing the coefficients of the basis elements $e_\mu \otimes e_\lambda$ on both sides, we obtain $P_\mu(P_\lambda(v)) = \delta_{\lambda, \mu} P_\lambda(v)$.
Using the left-to-right composition convention, this implies $P_\lambda \comp P_\mu = \delta_{\lambda, \mu} P_\lambda$.
This implies that the maps $\{P_\lambda\}_{\lambda \in \Lambda}$ form a family of mutually orthogonal projection operators summing to the identity.
Thus, $V$ completely decomposes into the direct sum of the images $\mathrm{Im}(P_\lambda)$.
It is straightforward to verify that $\mathrm{Im}(P_\lambda) = V_\lambda$, concluding the proof.
\end{proof}



\begin{defn}[Diagonalizable group scheme]
An affine group scheme $G$ over $\field$ is called \emph{diagonalizable} if there exists an abelian group $\Lambda$ such that
$G\cong\Spec(\field[\Lambda])$
as a group scheme. In this case, the group $\cha$ is canonically identified with the character group $\widehat{G}$ of $G$, and its coordinate ring $\mathcal{O}_G(G)$ is isomorphic to the group algebra $\field[\widehat{G}]$ as a Hopf algebra.
\end{defn}


\begin{lemma}[Weight space decomposition of global sections]
% ---- 原文行 2480 ----
\label{lem:GlobalSectionDecomposition}
Let $G$ be a diagonalizable group scheme acting on a scheme $X$, and let $(\mathcal{F}, \theta) \in \Mod_G(X)$ be a $G$-equivariant $\mathcal{O}_X$-module.
Then the space of global sections $V := \mathcal{F}(X)$($=\Gamma(X, \mathcal{F})$) naturally admits a right comodule structure over $\mathcal{O}_G(G)$, and therefore admits a canonical weight space decomposition
$$
\mathcal{F}(X) = \bigoplus_{\cha \in \widehat{G}}
\mathcal{F}(X)_{\cha}.
$$
where
$
\mathcal{F}(X)_{\cha}
:=
\left\{
s \in \mathcal{F}(X)
\left|
(\theta(G\times X))(\sigma^* s)
=
\left(\pr_G^{G \times X}\right)^*\cha
\cdot
\left(\pr_X^{G \times X}\right)^* s
%\in
% ---- 原文行 2500 ----
%\left(\left(\pr_X^{G \times X}\right)^* \mathcal{F}\right)(G\times X)
\right.\right\}.
$

\end{lemma}

\begin{proof}
The $G$-linearization $\theta : \sigma^* \mathcal{F} \xrightarrow{\cong} \left(\pr_X^{G \times X}\right)^* \mathcal{F}$ induces a natural map on the global sections.
For any section $s \in V = \mathcal{F}(X)$, its pullback along the action map $\sigma:G\times X \rightarrow X$ yields a section $\sigma^* s \in (\sigma^*\mathcal{F})(G\times X)$.
Applying the isomorphism $\theta$ to this pullback gives
$$
(\theta(G\times X))(\sigma^* s)
\in
\left(\left(\pr_X^{G \times X}\right)^* \mathcal{F}\right)(G\times X)
\cong
\mathcal{F}(X) \underset{\field}{\otimes} \mathcal{O}_G(G)
=
V \underset{\field}{\otimes} \mathcal{O}_G(G).
$$
This defines a $\field$-linear map
% ---- 原文行 2520 ----
$$
\rho
:
V \rightarrow V \underset{\field}{\otimes} \mathcal{O}_G(G),
\quad
s
\mapsto
(\theta(G\times X))(\sigma^* s).
$$


The cocycle condition and the unit condition of $\theta$
translates precisely to the coassociativity axiom
and the counit axiom of $\rho$,
respectively.
Therefore, $\rho$ defines a right comodule structure on $V$ over $\mathcal{O}_G(G) \cong \field[\widehat{G}]$.
By Proposition \ref{prop:WeightSpaceDecompositionForComodule},
this comodule structure completely decomposes $V$ into the direct sum of its weight spaces indexed by the character group $\widehat{G}$.

Under the canonical identification
% ---- 原文行 2540 ----
$
\left(\left(\pr_X^{G \times X}\right)^* \mathcal{F}\right)(G\times X)
\cong
\mathcal{F}(X) \underset{\field}{\otimes} \mathcal{O}_G(G),
$
the pure tensor $s \otimes e_{\cha} \in V \underset{\field}{\otimes} \mathcal{O}_G(G)$ corresponds to
$
\left(\pr_G^{G \times X}\right)^*\cha
\cdot
\left(\pr_X^{G \times X}\right)^* s
\in
\left(\left(\pr_X^{G \times X}\right)^* \mathcal{F}\right)(G\times X),
$
where
$\cdot$ denotes the scalar multiplication of elements of the $\mathcal{O}_{G\times X}(G\times U)$-module
$\left(\left(\pr_X^{G \times X}\right)^*\mathcal{F}\right)(G\times U)$
by elements of the ring $\mathcal{O}_{G\times X}(G\times U)$.
Thus, the condition $\rho(s) = s \otimes e_{\cha}$ translates into the condition
$$
(\theta(G\times X))(\sigma^* s) = \left(\pr_G^{G \times X}\right)^*\cha \cdot \left(\pr_X^{G \times X}\right)^* s,
% ---- 原文行 2560 ----
$$
establishing the formula for the weight space $\mathcal{F}(X)_{\cha}$.
\end{proof}





\begin{prop}[Weight space decomposition of hom spaces]
\label{prop:WeightSpaceDecompositionOfHomSpace}
Let $G$ be a diagonalizable group scheme acting on a scheme $X$.
For any
$(\mathcal{F}, \theta)$,
$(\mathcal{F}', \theta') \in \Mod_G(X)$
such that
$\mathcal{F}$
is finitely presented,
the hom space from $\mathcal{F}$ to $\mathcal{F}'$
in $\Mod(X)$
admits a canonical weight space decomposition
% ---- 原文行 2580 ----
$$
\operatorname{Hom}_{\Mod(X)}(\mathcal{F}, \mathcal{F}') 
= 
\bigoplus_{\cha \in \widehat{G}} \operatorname{Hom}_{\Mod_G(X)}\left( (\mathcal{F}, \theta)(\cha), (\mathcal{F}', \theta') \right).
$$
\end{prop}
\begin{proof}
Since $\mathcal{F}$ is finitely presented,
by Definition \ref{defn:InternalHomOfEquivariantModules}, 
the internal hom $\mathcal{H} := \mathcal{H}om_{\mathcal{O}_X}(\mathcal{F}, \mathcal{F}')$ is naturally equipped with a $G$-linearization $\Theta$ given by $\alpha \mapsto \theta^{-1} \comp \alpha \comp \theta'$, making $(\mathcal{H}, \Theta) \in \Mod_G(X)$.

By Lemma \ref{lem:GlobalSectionDecomposition},
the space of global sections
$$
\operatorname{Hom}_{\Mod(X)}(\mathcal{F}, \mathcal{F}')
=
\mathcal{H}(X)
$$
admits a canonical weight space decomposition
$
% ---- 原文行 2600 ----
\mathcal{H}(X)
=
\bigoplus_{\cha \in \widehat{G}} \mathcal{H}(X)_{\cha},
$
where
the following hold
for $f\in\mathcal{H}(X)$ and $\cha\in\widehat{G}$:
\begin{align*}
f\in \mathcal{H}(X)_{\cha}
\quad
&\Leftrightarrow
\quad
(\Theta(G\times X))(\sigma^* f)
=
\left(\pr_G^{G \times X}\right)^*\cha
\cdot
\left(\pr_X^{G \times X}\right)^* f
\\[2mm]
&\Leftrightarrow
\quad
% ---- 原文行 2620 ----
\theta^{-1} \comp \sigma^* f \comp \theta' = \left(\pr_G^{G \times X}\right)^*\cha \cdot \left(\pr_X^{G \times X}\right)^* f
\\[2mm]
&\Leftrightarrow
\quad
\sigma^* f \comp \theta' = \theta \comp \left( \left(\pr_G^{G \times X}\right)^*\cha \cdot \left(\pr_X^{G \times X}\right)^* f \right)
\\[2mm]
&\Leftrightarrow
\quad
\sigma^* f \comp \theta' = \left( \left(\pr_G^{G \times X}\right)^*\cha \cdot \theta \right) \comp \left(\pr_X^{G \times X}\right)^* f\\[2mm]
&(\because \left(\pr_G^{G \times X}\right)^*\cha \cdot\text{ is a scalar multiplication.})
\\[2mm]
&\Leftrightarrow
\quad
\sigma^* f \comp \theta' = \theta(\cha) \comp \left(\pr_X^{G \times X}\right)^* f
\quad(\because \text{definition of $\theta(\cha)$})
\end{align*}

This equality is precisely the condition for $f$ to be a $G$-equivariant morphism from
$(\mathcal{F}, \theta(\cha))$
to
% ---- 原文行 2640 ----
$(\mathcal{F}', \theta')$
(Definition \ref{defn:EquivariantSheaves}). 
Therefore, $\mathcal{H}(X)_{\cha} = \operatorname{Hom}_{\Mod_G(X)}\left( (\mathcal{F}, \theta)(\cha), (\mathcal{F}', \theta') \right)$, completing the canonical weight space decomposition.
\end{proof}
Finally, as a Corollary to the computation, we get an important intermediate step toward our proof of the main theorem.
\begin{cor}[Embedding of orbit categories]
\label{thm:EmbeddingOfTheOrbitCategory}
Let $G$ be a diagonalizable group scheme acting on a scheme $X$.
Then the forgetful functor $\Mod_G(X) \rightarrow \Mod(X)$ sending objects $(\mathcal{F}, \theta) \mapsto \mathcal{F}$ induces a fully faithful embedding of the orbit categories\footnote{
The same holds true for $\Coh_G(X)$ and $\Coh(X)$ when $X$ is locally Noetherian.
}:
$$
\QCoh_G(X)/\widehat{G}
\hookrightarrow
\QCoh(X).
$$
\end{cor}

\begin{proof}
By definition, the hom space in the orbit category is given by
% ---- 原文行 2660 ----
$$
\operatorname{Hom}_{\QCoh_G(X)/\widehat{G}}\left( (\mathcal{F}, \theta), (\mathcal{F}', \theta') \right)
=
\bigoplus_{\cha \in \widehat{G}} \operatorname{Hom}_{\QCoh_G(X)}\left( (\mathcal{F}, \theta)(\cha), (\mathcal{F}', \theta') \right).
$$
By Proposition \ref{prop:WeightSpaceDecompositionOfHomSpace}, this direct sum exactly coincides with the hom space $\operatorname{Hom}_{\QCoh(X)}(\mathcal{F}, \mathcal{F}')$ in $\QCoh(X)\subseteq\Mod(X)$.
Therefore, the functor $\QCoh_G(X)/\widehat{G} \hookrightarrow \QCoh(X)$ defined by $(\mathcal{F}, \theta) \mapsto \mathcal{F}$ on objects and by the identity map on morphisms is fully faithful.
\end{proof}



